\section{Topological Sort using DFS}
\label{sec:graph-topo-sort-dfs}

\problemstatement{Given a Directed Acyclic Graph (DAG) with $n$ nodes,
return a valid topological ordering -- a linear arrangement where
every edge $u \to v$ has $u$ appearing before $v$.}

\begin{patternbox}
\emph{``Order nodes so every prerequisite comes first''} is solved by
DFS with one addition: push a node onto a \textbf{stack} the moment it
\textbf{finishes} (every neighbor it can reach has already been fully
explored). Popping the stack from top to bottom yields a valid
topological order.
\end{patternbox}

\subsection*{Why finish-time ordering works}

When DFS finishes exploring node $u$ entirely (including everything
reachable from it) before pushing $u$ onto the stack, every node $v$
reachable from $u$ has \emph{already} been pushed (it must finish
before $u$ does, since $u$'s own finish waits on its recursive calls
returning). So $v$ ends up \textbf{lower} in the stack than $u$ -- and
popping top-to-bottom visits $u$ before $v$, exactly the required
order for edge $u \to v$.

\subsection*{Algorithm}

\begin{enumerate}
  \item Maintain a $\textit{visited}$ array and an empty stack.
  \item DFS from every unvisited node: recurse into all unvisited
        neighbors first, and only \textbf{after} that recursion
        returns (the node is now fully finished), push the current
        node onto the stack.
  \item Pop the entire stack to produce the topological order.
\end{enumerate}

\subsection*{C++ implementation}

\begin{lstlisting}
#include <bits/stdc++.h>
using namespace std;

void dfs(int node, vector<vector<int>>& adj, vector<bool>& visited, stack<int>& finishOrder) {
    visited[node] = true;
    for (int neighbor : adj[node]) {
        if (!visited[neighbor]) {
            dfs(neighbor, adj, visited, finishOrder);
        }
    }
    finishOrder.push(node); // push only after all neighbors are done
}

vector<int> topoSortDFS(int n, vector<vector<int>>& adj) {
    vector<bool> visited(n, false);
    stack<int> finishOrder;

    for (int i = 0; i < n; i++) {
        if (!visited[i]) {
            dfs(i, adj, visited, finishOrder);
        }
    }

    vector<int> result;
    while (!finishOrder.empty()) {
        result.push_back(finishOrder.top());
        finishOrder.pop();
    }
    return result;
}

int main() {
    int n = 6;
    vector<vector<int>> adj(n);
    adj[5] = {2, 0}; adj[4] = {0, 1}; adj[2] = {3}; adj[3] = {1};

    for (int x : topoSortDFS(n, adj)) cout << x << " ";
    cout << endl; // 5 4 2 3 1 0  (one valid order)
    return 0;
}
\end{lstlisting}

\subsection*{Dry run}

\dryrun{Take edges $5{\to}2$, $5{\to}0$, $4{\to}0$, $4{\to}1$,
$2{\to}3$, $3{\to}1$.}

DFS from $0$: no outgoing edges, finishes immediately, push $0$.
Stack: $[0]$. DFS from $1$: no outgoing edges, push $1$. Stack:
$[0,1]$. DFS from $2$: recurse to $3$; $3$ recurses to $1$ (already
visited, skip); $3$ finishes, push $3$. Stack: $[0,1,3]$. Back at $2$:
finishes, push $2$. Stack: $[0,1,3,2]$. DFS from $4$: neighbors $0,1$
both visited; finishes, push $4$. Stack: $[0,1,3,2,4]$. DFS from $5$:
neighbors $2,0$ both visited; finishes, push $5$. Stack:
$[0,1,3,2,4,5]$. Popping top to bottom: $5,4,2,3,1,0$ -- every edge
($5{\to}2$, $5{\to}0$, $4{\to}0$, $4{\to}1$, $2{\to}3$, $3{\to}1$)
indeed has its source appearing before its destination.

\begin{complexitybox}
\textbf{Time Complexity.} $O(V+E)$ -- standard DFS traversal cost.
\textbf{Space Complexity.} $O(V)$ for \textit{visited}, the stack, and
the recursion stack.
\end{complexitybox}

\begin{takeawaybox}
``Push on finish, read the stack top-down'' is a clean, purely local
rule -- no node needs any global information about the rest of the
graph to know when to push itself, only that its own DFS call has
returned. This locality is what makes the DFS approach to topological
sort so simple to implement correctly.
\end{takeawaybox}
